% GENERATED FILE. Edit canonical lessons, then run npm run generate:books.
% canonical-source-hash: fnv1a64:3062c05174f48c62
% canonical-lessons: TE-C63-hunger, TE-C63-sleep, TE-C63-cough, TE-C63-disease, TE-C63-health

\chapter{Hunger, Sleep and Health}
\label{ch:te-health}
\hlchaptermodality{63}

\begin{chapteropening}
\textbf{By the end of this chapter:} \emph{I can say in Telugu that I am hungry, sleepy or unwell, and ask after somebody else's health.}

Complete TE-C63-health and demonstrate the chapter promise.
\end{chapteropening}

\section[ākali]{\te{ఆకలి} (ākali) — hunger}

\label{lesson:TE-C63-hunger}



\begin{quote}
Before the new one: what did \emph{addaṁ} mean?
\end{quote}

\subsection*{You'll want to know: \te{ఆకలి}}
\textbf{\te{ఆకలి}} (\emph{ākali}) — \textquotedblleft{}hunger\textquotedblright{}.

\te{ఆకలి} is hunger. Telugu does not say \textquotedblleft{}I am hungry\textquotedblright{} — it says \te{నాకు} \te{ఆకలి}, \textquotedblleft{}to me, hunger\textquotedblright{}, the same shape as \te{నాకు} \te{తెలుసు} and \te{నాకు} \te{తెలుగు} \te{వచ్చు}. Hunger is something that happens to you.

The word is inherited, and it is among the first a child in a Telugu house learns to say — usually loudly, and usually before \emph{bhōjanaṁ} is anywhere near ready.

\te{ఆకలిగా} \te{ఉందా}? is a real greeting in a house you are close to. It gets asked before \te{మీరు} \te{ఎలా} \te{ఉన్నారు}? has a chance.

\begin{grammarlens}[title={What you've built}]
The first thing to say about yourself.
\end{grammarlens}

\subsection*{Your turn}
Say these aloud:

\begin{itemize}
\raggedright
  \item \emph{ākali}
  \item it once more, slowly
  \item \emph{nāku ākali}, then ask \emph{mīru elā unnāru?}
\end{itemize}

\begin{tcolorbox}[breakable,colback=teal!4,colframe=teal!35!black,title={Before you move on}]
What does \te{ఆకలి} mean? (\textquotedblleft{}hunger\textquotedblright{}.) Say it once more.
\end{tcolorbox}
\section[nidra]{\te{నిద్ర} (nidra) — sleep}

\label{lesson:TE-C63-sleep}



\begin{quote}
Before the new one: what did \emph{ākali} mean?
\end{quote}

\subsection*{You'll want to know: \te{నిద్ర}}
\textbf{\te{నిద్ర}} (\emph{nidra}) — \textquotedblleft{}sleep\textquotedblright{}.

\te{నిద్ర} is sleep, and it takes the same shape as \te{ఆకలి}: \te{నాకు} \te{నిద్ర} \te{వస్తోంది}, \textquotedblleft{}sleep is coming to me\textquotedblright{}. You do not do it; it arrives.

The word came from Sanskrit \emph{nidrā} and settled without changing. Beside it Telugu keeps \te{నిదుర}, the softer everyday form, and the two are one word wearing different clothes.

\te{నిద్ర} also makes a polite reason for leaving late in the evening, and it works the way \te{గాలి} worked earlier. Nobody argues with it.

\begin{grammarlens}[title={What you've built}]
Two, and both of them come to you.
\end{grammarlens}

\subsection*{Your turn}
Say these aloud:

\begin{itemize}
\raggedright
  \item \emph{nidra}
  \item it once more, slowly
  \item \emph{nidra}, then \emph{ākali}, and say which one you have now
\end{itemize}

\begin{tcolorbox}[breakable,colback=teal!4,colframe=teal!35!black,title={Before you move on}]
What does \te{నిద్ర} mean? (\textquotedblleft{}sleep\textquotedblright{}.) Say it once more.
\end{tcolorbox}
\section[daggu]{\te{దగ్గు} (daggu) — a cough}

\label{lesson:TE-C63-cough}



\begin{quote}
Before the new one: what did \emph{nidra} mean?
\end{quote}

\subsection*{You'll want to know: \te{దగ్గు}}
\textbf{\te{దగ్గు}} (\emph{daggu}) — \textquotedblleft{}a cough\textquotedblright{}.

\te{దగ్గు} is a cough. It is inherited, and it is a sound turned into a word, which is how a good many of these words start out.

Like \te{ఆకలి} and \te{నిద్ర} it comes to you rather than from you: \te{నాకు} \te{దగ్గు}, and nobody needs telling any more than that.

Telugu keeps \te{దగ్గు} for the whole thing at once — the single cough, the fit of coughing, and the illness sitting behind both — where English reaches for three phrases. A \te{మంచు} morning in the hills fills a house with it.

\begin{grammarlens}[title={What you've built}]
Three.
\end{grammarlens}

\subsection*{Your turn}
Say these aloud:

\begin{itemize}
\raggedright
  \item \emph{daggu}
  \item it once more, slowly
  \item \emph{daggu}, then \emph{nidra}, and say which one keeps you from the other
\end{itemize}

\begin{tcolorbox}[breakable,colback=teal!4,colframe=teal!35!black,title={Before you move on}]
What does \te{దగ్గు} mean? (\textquotedblleft{}a cough\textquotedblright{}.) Say it once more.
\end{tcolorbox}
\section[rōgaṁ]{\te{రోగం} (rōgaṁ) — a disease}

\label{lesson:TE-C63-disease}



\begin{quote}
Before the new one: what did \emph{daggu} mean?
\end{quote}

\subsection*{You'll want to know: \te{రోగం}}
\textbf{\te{రోగం}} (\emph{rōgaṁ}) — \textquotedblleft{}a disease\textquotedblright{}.

\te{రోగం} is illness in general: not the \te{దగ్గు} you have this week, but the thing a \te{వైద్యుడు} puts a name to.

It is Sanskrit, from a root about breaking — an illness is something breaking in you — and Telugu took the word whole, the way it took \emph{hṛdayaṁ}.

\te{రోగి} is the person who has one. A \te{వైద్యుడు}'s whole working day is spent between those two words, and so is a good deal of ordinary conversation about relatives.

\begin{grammarlens}[title={What you've built}]
Four.
\end{grammarlens}

\subsection*{Your turn}
Say these aloud:

\begin{itemize}
\raggedright
  \item \emph{rōgaṁ}
  \item it once more, slowly
  \item \emph{rōgaṁ}, then \emph{daggu}, and say which one is the smaller trouble
\end{itemize}

\begin{tcolorbox}[breakable,colback=teal!4,colframe=teal!35!black,title={Before you move on}]
What does \te{రోగం} mean? (\textquotedblleft{}a disease\textquotedblright{}.) Say it once more.
\end{tcolorbox}
\section[ārōgyaṁ]{\te{ఆరోగ్యం} (ārōgyaṁ) — health}

\label{lesson:TE-C63-health}



\begin{quote}
Before the new one: what did \emph{rōgaṁ} mean?
\end{quote}

\subsection*{You'll want to know: \te{ఆరోగ్యం}}
\textbf{\te{ఆరోగ్యం}} (\emph{ārōgyaṁ}) — \textquotedblleft{}health\textquotedblright{}.

\te{ఆరోగ్యం} is health, and it is the word before it with a small piece stuck on the front. The piece means \textquotedblleft{}not\textquotedblright{}, and \te{రోగ} is the illness: health here is what is left when the illness is gone.

That is a real difference in outlook and worth a moment. English health is something to be built up. \te{ఆరోగ్యం} is a subtraction.

\te{మీ} \te{ఆరోగ్యం} \te{ఎలా} \te{ఉంది}? is what a house asks an \emph{atithi} who has been ill. It is more particular than \te{మీరు} \te{ఎలా} \te{ఉన్నారు}?, and it expects a real answer.

\begin{grammarlens}[title={What you've built}]
Five: \te{ఆకలి}, \te{నిద్ర}, \te{దగ్గు}, \te{రోగం}, \te{ఆరోగ్యం}. Enough to say what is wrong and to ask after somebody else.
\end{grammarlens}

\subsection*{Your turn}
Say these aloud:

\begin{itemize}
\raggedright
  \item \emph{ārōgyaṁ}
  \item it once more, slowly
  \item all five in order, then ask \emph{mī ārōgyaṁ elā undi?}
\end{itemize}

\begin{tcolorbox}[breakable,colback=teal!4,colframe=teal!35!black,title={Before you move on}]
What does \te{ఆరోగ్యం} mean? (\textquotedblleft{}health\textquotedblright{}.) Say it once more.
\end{tcolorbox}
